\scanpage{31}
2. \uline{Familles de classes de cycles \deleted{\unsure{à équiv.}} \inserted{mod équiv. d'Alb.}}

\uline{Théorème} 2.1. Soient \(X^n\) et \(Y^m\) comme d'habitude, et \(Z\in\deleted{\unsure{\operatorname{Cy}}}A^i_{\inserted{\text{\textcircled{alb}}}}(X\times Y)\), d'où une \deleted{homomorphisme}\inserted{application} \(\nu\longmapsto Z(\nu)\) de \(\deleted{\unsure{A_0(Y)}}\inserted{A_{0\,\mathrm{alb}}(Y)}\) dans \(A^i_{\mathrm{alb}}(X)\).
\deleted{Il existe une application et une seule de \((A^0_{\mathrm{alb}}(Y)\simeq\operatorname{Alb}_Y^*(k))\) dans \(A^i_{\mathrm{alb}}(X)\) telle que pour tout \(\nu\in\unsure{\operatorname{Alg}}_0(Y)\)}
Or \(A_{0\,\mathrm{alb}}(Y)\overset{S_Y}{\simeq}\deleted{\unsure{B}}A_Y^*(k)\), d'où un hom de groupes
\diagram{\begin{tikzcd}[ampersand replacement=\&,execute at end picture={\draw ([xshift=3mm]ay.west) ellipse (3mm and 4mm);\node at ([xshift=3mm,yshift=-6mm]ay.west) {Alb};}]
|[alias=ay,label={left:$u_Z^*:\,$}]| A_Y^*(k)\arrow[r]\&A^i_{\mathrm{alb}}(X)
\end{tikzcd}}
\deleted{dont la restriction à \(A_Y^{(n)}(k)\) est} défini pr la condition
\[
\deleted{\unsure{S_Y}}Z(\nu)=u_Z(S(\nu))
\]
pr tt \(\nu\) (resp. pr tt \(\nu\) de degré \(n\), si on regarde la restriction en degré \(n\)).

\uline{De plus}, si on regarde l'hom \uline{induit} sur \(A_Y(k)=A_Y^0(k)\)
\[
u_Z:A_Y(k)\longrightarrow A^i_{\mathrm{alb}}(X),
\]
on a ce qui suit : \deleted{Il existe un hom de passage de VA \(\varphi:A'\longrightarrow A_Y\) et un cycle \(Z'\in A^i_{\mathrm{alb}}(\,\unsure{/}\)}

(i) Il ne dépend que de la classe de \(Z\) dans \(A^{i\,\unsure{\tau}}_{\deleted{\unsure{\mathrm{alg}}}}(X\times Y)\)
\margin{[À \unsure{réserver} pr plus tard}

(ii) Il existe une \deleted{hom}\inserted{isogénie}
\mspage{30}
